\chapter{Fejlesztői dokumentáció}
\label{ch:dev_doc}

Ebben a fejezetben az alkalmazás technikai megvalósításának részleteit mutatom be.


\subsection{Célközönség és a szoftver célkitűzései}
\label{subsec:target_audience}

Az alkalmazás elsődleges célja, hogy modern, digitális alternatívát nyújtson a snapszer kártyajátéknak. Lényegében azoknak készült, akik szeretnének akkor is játszani, amikor nincsenek közvetlenül körülöttük azok a személyek (például családtagok vagy barátok), akikkel hagyományosan, fizikai kártyapaklival szoktak játszani. Az alkalmazás áthidalja a fizikai távolságot, és lehetőséget ad a közös kikapcsolódásra akkor is, ha a felek nincsenek egy helyiségben. A fejlesztés során négy fő célcsoportot határoztam meg, akiknek az alkalmazás funkcionalitását és felületét dedikáltam:
\begin{itemize}
    \item \textbf{A klasszikus kártyajátékok kedvelői:} Azok a felhasználók, akik ismerik és szeretik a snapszer szabályait, de a mindennapok során nehézséget okoz számukra fizikai partnert találni a játékhoz. A snapszer.com élesített rendszere nekik biztosít azonnali párosítási lehetőséget a várólista segítségével.
    
    \item \textbf{Digitális platformot kereső közösségek:} Az alkalmazás célközönségébe tartoznak azok a baráti körök vagy ismerősök, akik földrajzi távolságtól függetlenül szeretnének egymás ellen játszani. A rendszer regisztrációs és statisztikai modulja lehetővé teszi számukra a hosszú távú versengést és az eredmények visszakövetését.
    
    \item \textbf{Alacsony belépési küszöböt igénylő felhasználók:} Mivel az applikáció böngésző alapú és nem igényel telepítést, célközönségnek tekinthetők a kevésbé technikai beállítottságú játékosok is.
    
    \item \textbf{A játékot tanulni vágyók (kezdők):} A felület ideális azoknak is, akik most szeretnék megtanulni a snapszer alapjait. Az alkalmazás rendelkezik szabályzat oldallal, amely segít az elméleti tudás elsajátításában, a szoftverbe épített automatizált szabályellenőrzés pedig a gyakorlatban tanítja meg az érvényes lépéseket, megakadályozva a szabálytalan kártyarakást.
    
\end{itemize}

A szoftver végső célja egy olyan közösségi tér létrehozása, ahol a hagyományőrzés és a modern technológia találkozik, biztosítva a játék tisztaságát az automatizált szabályellenőrzéssel és az igazságos pontszámítással. 


\section{Implementációt megelőző rendszertervek}
\label{sec:pre_impl_design}

Ez a szakasz azokat az előzetes terveket tartalmazza, amelyek a kódolás megkezdése előtt készültek, és referenciaként szolgáltak a fejlesztés során. Az ábrák a tervezési fázisból és az szoftver architektúrája fázisból készültek.

\subsection{Használatieset-diagram (Use-case diagram)}

Ez a diagram a fejlesztés megkezdése előtt rögzített funkcionális elvárásokat és a felhasználói szerepkörök interakcióit szemlélteti.

\begin{figure}[H]
    \centering
    \includegraphics[width=0.6\textwidth]{use_case_pre}
    \caption{Előzetes használatieset-diagram}
    \label{fig:pre_use_case}
    \centering{\footnotesize Az ábra a draw.io \cite{drawio} alkalmazás segítségével készült.}
\end{figure}

A használatieset-diagram (Use-case diagram) funkcionális leírását az alábbi táblázat egészíti ki:

\begin{table}[H]
  \centering
  \begin{tabular}{ | m{0.28\textwidth} | m{0.32\textwidth} | m{0.30\textwidth} | }
    \hline
    \textbf{Given (Adott)} & \textbf{When (Amikor)} & \textbf{Then (Akkor)} \\
    \hline \hline
    Vendégként tartózkodom a regisztrációs oldalon. & Kitöltöm a szükséges adatokat (felhasználónév, érvényes e-mail cím, jelszó) és rákattintok a „Regisztráció” gombra. & A Rendszer létrehozza a fiókomat. \\
    \hline
    Vendégként tartózkodom a bejelentkezési oldalon. & Megadom a regisztrált felhasználónevemet és jelszavamat, és rákattintok a „Belépés” gombra. & A Rendszer hitelesíti az adataimat, és sikeresen a főmenübe irányít. \\
    \hline
    Játékosként be vagyok jelentkezve, és a főmenüben vagyok. & Rákattintok a játék keresése gombra. & A Rendszer létrehoz egy szobát és keres nekem ellenfelet. \\
    \hline
    Játékosként egy játékban vagyok, és én vagyok soron. & A kezemben lévő kártyák közül érvényesen kijátszok egyet. & A Rendszer elfogadja a lépést, a kijátszott kártya megjelenik az asztalon, és a sor a következő játékosra száll. \\
    \hline
    Játékosként van egy királyom és egy felsőm azonos színből és én vagyok soron. & Egy gombot kattintva kijátszom a 20/40-et és kiválasztom a királyt vagy a felsőt. & A Rendszer elfogadja a lépést, hozzárendeli a pontokat és megjeleníti a lapot az asztalon. \\
    \hline
    Két játékos kijátszotta a kártyáját az adott körben. & A Rendszer kiértékeli a két kijátszott kártyát a snapszer szabályai szerint. & A Rendszer megállapítja a kör nyertesét, kiosztja a pontokat, és frissíti a kijelzőt. \\
    \hline
    Véget ér egy kör. & A kört nyert játékosnak összegyűlik a 66 pont. & A rendszer kihirdeti a játékost, aki nyerte a partit. \\
    \hline
    Játékosként egy játékteremben tartózkodom, de még nem fejeződött be a parti. & Rákattintok a „Játék elhagyása” gombra. & A Rendszer eltávolít a játékteremből és átirányít a főmenübe. \\
    \hline
    Játékosként bejelentkezve a főmenüben tartózkodom. & Rákattintok a statisztikák gombra. & Megjelenik egy oldal, amin az eddig elért partikból láthatóak a statisztikák. \\
    \hline
  \end{tabular}
  \caption{A rendszer használati eseteinek részletezése}
  \label{tab:use-case-scenarios}
  \centering{\footnotesize A táblázat a Microsoft Word \cite{msword} alkalmazás segítségével készült.}
\end{table}

\subsection{A felhasználói felület tervei (Figma)}

A vizuális tervezési szakaszban készült drótváztervek az alkalmazás ergonómiáját és a tervezett komponensek elrendezését rögzítették. Az összes drótváz ábrát a Figma \cite{figma} grafikai tervező szoftver segítségével készítettem el.

\vspace{1em}
\noindent\textbf{Bejelentkezés nélküli főoldal}

\begin{figure}[H]
    \centering
    \includegraphics[width=0.8\textwidth]{bejelentkezes_nelkuli_figma}
    \caption{Bejelentkezés nélküli főoldal drótváz ábrája}
    \label{nologinfigma}
    \vspace{0.5em}
\end{figure}

\vspace{1em}
\noindent\textbf{Bejelentkezett főoldal}

\begin{figure}[H]
    \centering
    \includegraphics[width=0.8\textwidth]{bejelentkezett_menu_figma}
    \caption{Bejelentkezett főoldal drótváz ábrája}
    \label{loggedinfigma}
    \vspace{0.5em}
\end{figure}

\vspace{6em}
\noindent\textbf{Bejelentkezés oldal}

\begin{figure}[H]
    \centering
    \includegraphics[width=0.8\textwidth]{bejelentkezes_figma}
    \caption{Bejelentkezés oldal drótváz ábrája}
    \label{loginpage}
    \vspace{0.5em}
\end{figure}

\vspace{1em}
\noindent\textbf{Regisztráció oldal}

\begin{figure}[H]
    \centering
    \includegraphics[width=0.8\textwidth]{regisztracio_figma}
    \caption{Regisztráció oldal drótváz ábrája}
    \label{registerpage}
    \vspace{0.5em}
\end{figure}

\vspace{20em}
\noindent\textbf{Játék oldal}

\begin{figure}[H]
    \centering
    \includegraphics[width=0.8\textwidth]{jatekfelulet}
    \caption{Játékoldal drótváz ábrája, amikor az ellenfél van soron}
    \label{gamepage_opp_turn}
    \vspace{0.5em}
\end{figure}

\begin{figure}[H]
    \centering
    \includegraphics[width=0.8\textwidth]{jatekfelulet_mi_jovunk}
    \caption{Játékoldal drótváz ábrája, amikor mi vagyunk soron}
    \label{gamepage_our_turn}
    \vspace{0.5em}
\end{figure}

\subsection{Komponens diagramok}

\vspace{1em}
\noindent\textbf{Magas szintű architektúra} \\
Az első diagram a tervezett felhasználói felületet (Client/Frontend), a kiszolgáló egységet (Backend) és az adatbázist (DB) és a kapcsolatait szemlélteti.

\begin{figure}[H]
    \centering
    \includegraphics[width=0.85\textwidth]{komponens_diagram_pre}
    \caption{A rendszer három fő komponensének kapcsolata}
    \label{pre_comp_1}
    \vspace{0.5em}
    \centering{\footnotesize Az ábra a draw.io \cite{drawio} alkalmazás segítségével készült.}
\end{figure}

\vspace{1em}
\noindent\textbf{A Backend részletes felépítése} \\
A második diagram részletes betekintést nyújt a tervezett Backend belső működésébe, elmagyarázva az API vezérlőt és a valós idejű kommunikációért felelős Socket modul feladatait.

\begin{figure}[H]
    \centering
    \includegraphics[width=0.85\textwidth]{komponens_diagram_backend_pre}
    \caption{A Backend belső moduláris felépítésének terve}
    \label{pre_comp_2}
    \vspace{0.5em}
    \centering{\footnotesize Az ábra a draw.io \cite{drawio} alkalmazás segítségével készült.}
\end{figure}
\subsection{Osztálydiagram (UML)}

Az előzetes osztálydiagram a szoftver belső objektumorientált adatszerkezeteknek a tervét mutatja be.

\begin{figure}[H]
    \centering
    \includegraphics[width=0.9\textwidth]{uml_pre}
    \caption{Tervezett osztálydiagram (UML)}
    \label{fig:pre_uml_class}
    \centering{\footnotesize Az ábra a draw.io \cite{drawio} alkalmazás segítségével készült.}
\end{figure}

\subsection{Egyed-kapcsolat diagram (ER diagram)}

Az adatbázis-tervezés során készült ER diagram a tárolni kívánt adatok közötti logikai relációkat és a táblák tervezett sémáját ábrázolja.

\begin{figure}[H]
    \centering
    \includegraphics[width=1\textwidth]{ER_pre}
    \caption{Az adatbázis előzetes logikai terve (ER diagram)}
    \label{fig:pre_er_diagram}
    \vspace{0.5em}
    \centering{\footnotesize Az ábra a draw.io \cite{drawio} alkalmazás segítségével készült.}
\end{figure}

\section{A szoftver belső struktúrája és adatábrázolása}
\label{sec:internal_structure}

A rendszer stabilitását és teljesítményét nagymértékben meghatározza az alkalmazott szoftverarchitektúra, ezért ebben a szakaszban a backend forráskódjának felépítését mutatom be. A kódolás során a modularitás elve érvényesült, így az egyes komponensek könnyen tesztelhetők, és a jövőbeli továbbfejlesztések alkalmával is akadálymentesen bővíthetők.

\subsection{Moduláris felépítés és kódminőség}
\label{subsec:modularity}

Jelentősen megnőtt a kód olvashatósága és karbantarthatósága, mivel szigorúan szét lettek választva a logikai rétegek a szerveroldalon. Különálló, objektumorientált osztályokba (például \texttt{Card.js}, \texttt{Deck.js}, \texttt{Player.js}) került a játékmenet alapvető logikája, így ezek a fájlok belső állapottal (encapsulation) rendelkeznek, és teljesen függetlenek a hálózati keretrendszertől. Kimagasló újrafelhasználhatóságot biztosít ez a megközelítés, ugyanis a központi játékmotor (\texttt{Engine.js}) lényegében bármilyen környezetben (akár egy lokális, asztali alkalmazásban is) futtatható lenne kódmódosítás nélkül. Megkönnyítené egy mobilalkalmazás jövőbeli fejlesztését is ez az elszigetelt felépítés, mivel a játékmotor kódja könnyedén újrahasznosítható.

Kifejezetten biztonságos és robusztus a hitelesítési folyamat is, mivel az \texttt{authController.js} aszinkron vezérlőfunkciói felelnek a jelszavak hasheléséért és az adatok validálásáért. Garantált a szerveroldali adatintegritás, hiszen az adatbázisba történő írás során a kód a \texttt{sequelize.transaction} metódust alkalmazza, amely biztosítja, hogy szerverhiba esetén a félig végrehajtott módosítások automatikusan visszavonásra kerüljenek (rollback).

\subsection{Adatábrázolási döntések (Sequelize ORM)}
\label{subsec:data_representation}

Kizárólag relációs adatbáziskezelőt alkalmaz a rendszer háttértára, amelynek a dinamikus eléréséhez és az adattípusok absztrakt definiálásához a Sequelize ORM szoftverkönyvtárat implementáltam\cite{sequelize}. JavaScript osztályokként modelleztem az entitásokat, hogy a fejlesztés során ne kelljen nyers SQL parancsokat írni, ezáltal védem az alkalmazást az SQL injection támadás ellen.

Hatékonyan valósulnak meg a táblák közötti összeköttetések, mert az \texttt{index.js} fájlban pontosan lettek definiálva a relációk, mint amilyen a \texttt{Player} és a \texttt{History} (mérkőzés-történet) modellek közötti egy-a-többhöz (1:N) kapcsolat is. Soha nem maradnak árva rekordok az adatbázisban a profilok törlése után, mivel a kapcsolatoknál minden esetben be lett állítva az \texttt{ON DELETE CASCADE} szabály, amely a szülő entitás eltávolításakor a kapcsolódó al-adatokat is megsemmisíti. Modern és skálázható megoldásnak számít az elsődleges kulcsok (Primary Key) ábrázolása is, ugyanis a hagyományos inkrementális sorszámozás helyett automatikusan generált, prediktálhatatlan \texttt{UUIDv4} azonosítókat (\texttt{DataTypes.UUID}) kapnak a rekordok az adatbázisban.

\subsection{Kliensoldali architektúra és állapotkezelés}
\label{subsec:frontend_architecture}

Az alkalmazás a teljes útvonalválasztást a \texttt{React Router} keretrendszerel kezeli az \texttt{App.jsx} fájlban, oldalújratöltések nélkül. A frontend kódja moduláris, hiszen a felhasználói felület önálló, funkcióspecifikus React komponensekre (például \texttt{GamePage}, \texttt{StatsPage}, \texttt{GuestPage}) bomlik. 

A valós idejű hálózati kommunikációt a \texttt{socket.js}-ben inicializált Socket.IO \cite{socketio} kliens végzi, amelynek eseményeit a React beépített \texttt{useEffect} hookjai figyelik. Különösen megbízható a vizuális visszacsatolás a játékmenet során, mert a szabálytalan lépéseket jelző animációkhoz és az értesítésekhez tartozó időzítéseket (timeouts) a \texttt{useRef} hookokkal izoláltan, a memóriában tartva kezeli a rendszer, elkerülve a memóriaszivárgást.

A JWT tokenek a böngésző \texttt{localStorage} tárolójába kerülnek, így magas szintű biztonságot és konzisztens felhasználói élményt garantál a munkamenet-kezelés. Szinkronban maradnak a megnyitott böngészőlapok, mivel a fő \texttt{App} komponens egy globális eseményfigyelővel monitorozza a tárhely változásait, így egy esetleges kijelentkezés azonnal és automatikusan érvényesül a felhasználó összes nyitott ablakában. A megjelenés reszponzív és modern az összes nézeten, erre a Tailwind CSS \cite{tailwindcss} keretrendszert alkalmaztam, amit közvetlenül a JSX struktúrába integráltam.

\subsection{Fizikai komponens-telepítési diagram}
\label{subsec:deployment_diagram}

Jól áttekinthető képet nyújt az infrastruktúráról a fizikai telepítési diagram, amely megmutatja a Microsoft Azure felhőszolgáltatásokba kihelyezett szoftverkomponensek topológiáját.

\begin{figure}[H]
    \centering
    % TODODODODDOIde jöhet majd a PlantUML által generált telepítési diagram hivatkozása
    % \includegraphics[width=0.9\textwidth]{deployment_diagram}
    \caption{A rendszer fizikai telepítési diagramja (Deployment Diagram)}
    \label{fig:deployment_diagram}
\end{figure}

 Az alkalmazás könnyen skálázható, mivel a frontend statikus fájljait és a Node.js API szervert egy hosztolt App Service környezet szolgálja ki, a perzisztens adattárolásért pedig egy tőle fizikailag független, de azonos virtuális hálózaton elhelyezkedő Azure Database for PostgreSQL szerver felel.

\section{A megvalósult rendszerstruktúra}
\label{sec:actual_implementation}

A fejlesztési folyamat során az előzetes tervek bizonyos pontokon változtak és alkalmazkodtak a technológiai kihívásokhoz, valamint a felmerülő új kihívásokhoz. Az alábbi diagramok a szoftver végleges, implementált állapotát tükrözik.

\subsection{Végleges osztálydiagram (UML)}

A megvalósult osztálydiagram bemutatja a tényleges az üzleti logikát kiszolgáló osztályokat, a köztük lévő tényleges kapcsolatokat. A diagram a végleges kódbázis alapján készült.

\begin{figure}[H]
    \centering
    \includegraphics[width=0.9\textwidth]{UML}
    \caption{A rendszer megvalósult osztálydiagramja}
    \centering{\footnotesize Az ábra a PlantUML \cite{plantuml} eszköz segítségével készült.}
    \label{fig:actual_uml}
\end{figure}

Az osztálydiagramon szereplő komponensek a felelősségi körök egyértelmű szétválasztásának elve alapján lettek kialakítva. A rendszer főbb osztályai és azok funkciói a következők:

\subsection{Osztályok}

\begin{itemize}
    \item \textbf{Card:} Az osztályok közül a legkisebb, amely egyetlen kártyalapot reprezentál. Tárolja a kártya színét, típusát és a játékban betöltött pontértékét . Képes meghatározni önmagáról, hogy az aktuális aduszín ismeretében ő maga adunak minősül-e.
    
    \item \textbf{Deck:} A kártyapakli állapotának és műveleteinek kezeléséért felelős osztály. Feladata a snapszer szabályainak megfelelő 20 lapos csomag generálása, megkeverése, valamint a kártyahúzás (drawCard) és a pakli ürülésének (isEmpty) logikai biztosítása. A keveréshez a Knuth keverési algoritmust alkalmaztam.
    
    \begin{algorithm}[H]
    \caption{Fisher--Yates (Knuth) keverés algoritmus}
    \label{alg:fisher_yates}
    \begin{algorithmic}[1]
    \Function{Shuffle}{$A$ tömb}
        \State $n \gets \text{hossz}(A)$
        \For{$i \gets n - 1 \text{ downto } 1$}
            \State $j \gets \text{véletlen egész szám } 0 \text{ és } i \text{ között (inkluzív)}$
            \State Cseréljük meg az $A[i]$ és $A[j]$ elemeket
        \EndFor
    \EndFunction
    \end{algorithmic}
    \end{algorithm}
    \centering{\footnotesize Az algoritmus forrása a Wikipédia\cite{fisher_yates_wiki}.}
    
    \item \textbf{Player:} Egy adott mérkőzésben résztvevő kliens szerveroldali reprezentációja. Nyilvántartja a játékos kezében lévő lapokat, a már megnyert (ütött) kártyákat, valamint az aktuális körben és a meccs során szerzett pontokat. Ez az osztály végzi a játékoshoz köthető specifikus ellenőrzéseket is, például a 20-as és 40-es mondások (announcements) felismerését és validálását.
    
    \item \textbf{Engine:} A játékszerver központi magja, amely a teljes üzleti logikát és a klasszikus Snapszer játékszabályokat implementálja. Független a hálózati rétegtől, kizárólag a játék menetéért felel. Feladata a körök elindítása és lezárása, az ütés- és színkényszer szigorú betartatása (isValidMove), a kártyalapok erősségének összehasonlítása, a pakli takarása (close), valamint a pontszámítás és a győzelmi feltételek ellenőrzése.
    
    \item \textbf{Room:} A hálózati kommunikáció és a játékmotor (Engine) közötti híd szerepét tölti be. Minden elindult mérkőzéshez létrejön egy szobapéldány, amely a Socket.IO eseményeket kezeli. Fogadja a kliensek lépéseit (kártyarakás, aducsere), továbbítja azokat validálásra a játékmotornak, majd a frissített játékállapotot (broadcastState) elküldi minden játékeseménynél a két kliensnek.
\end{itemize}

\subsection{Végleges adatbázis-séma (ER diagram)}

Az alkalmazás adatperzisztenciáját biztosító adatbázis végleges szerkezete az alábbi egyed-kapcsolat diagramon (ER diagram) látható. Az ábra szemlélteti a létrehozott táblákat, a tárolt adattípusokat, az elsődleges és idegen kulcsokat, valamint a köztük lévő relációkat.

\begin{figure}[H]
    \centering
    \fbox{\includegraphics[width=1\textwidth]{ER}}
    \caption{A megvalósult adatbázis szerkezete (ER diagram)}
    \label{fig:actual_er}
    \vspace{0.5em}
    \centering{\footnotesize Az ábra a draw.io \cite{drawio} alkalmazás segítségével készült, a DBeaver \cite{dbeaver} által generált DDL utasítások felhasználásával.}
\end{figure}

\subsection{Az adatbázis táblái}

Az ER diagramon ábrázolt adatbázistáblák a relációs adatábrázolás és a felelősségi körök szétválasztásának elveit követve kerültek kialakításra. A rendszer entitásai a következők:


\begin{itemize}
    \item \textbf{auth:} A hitelesítési adatok és a felhasználói fiókok biztonságos tárolásáért felelős tábla. Tartalmazza a rekord egyedi azonosítóját (id), az egyedi (UNIQUE kényszerrel védett) felhasználónevet (username) és a bejelentkezéshez szükséges jelszót (password). Kizárólag az autentikációs folyamatok során van szerepe.
    
    \item \textbf{player:} A játékosok publikus profiladatait tartalmazó tábla. Minden bejegyzése pontosan egy hitelesítési fiókhoz kapcsolódik az \texttt{auth\_id} idegen kulcson keresztül (1:1 kapcsolat). A profil létrehozásának idején (created\_at) felül redundánsan tartalmazza a felhasználónevet is.
    
    \item \textbf{stats:} Az egyes játékosokhoz tartozó globális és összesített statisztikákat nyilvántartó entitás. A player\_id idegen kulcs segítségével rendelődik hozzá a megfelelő profilhoz. Célja a teljesítménykövetés: tárolja az eddig lejátszott mérkőzések számát (total\_games), a megszerzett győzelmeket (total\_wins) és az összesített meccspontokat (total\_points).
    
    \item \textbf{history:} A korábban lejátszott mérkőzések történetét és eredményeit rögzítő tábla, amely egy-a-többhöz (1:N) kapcsolatban áll a \texttt{player} táblával. Minden lezárt mérkőzés után új rekord jön létre benne, amely megőrzi a mérkőzés időpontját (match\_date), az aktuális ellenfél nevét (opponent\_name), a meccs végkimenetelét (match\_status), valamint az elért végeredményt (final\_score).
\end{itemize}

\section{Backend architektúra és infrastruktúra}

\subsection{Fontosabb metódusok dokumentációja}
\label{subsec:method_docs}

A játékmotor belső működésének megértéséhez elengedhetetlen a kulcsfontosságú metódusok pontos definíciója. Az alábbiakban a rendszer legfontosabb ellenőrző logikáját mutatom be.

\begin{table}[H]
    \centering
    \begin{tabular}{| p{0.22\textwidth} | p{0.73\textwidth} |}
        \hline
        \multicolumn{2}{|l|}{\textbf{Az \texttt{Engine} osztály \texttt{isValidMove} metódusának megvalósítása:}} \\
        \hline
        \textbf{Bemenő adat} & \texttt{player: Player} -- a lépést végrehajtó játékos objektuma; \texttt{cardToPlay: Card} -- a kijátszani kívánt kártyalap; \texttt{ledCard: Card} -- az ellenfél által hívott (már az asztalon lévő) kártyalap \\
        \hline
        \textbf{Kimenő adat} & \texttt{boolean} -- \texttt{true}, ha a kártya kijátszása a szabályoknak megfelel; \texttt{false}, ha az adott szituációban a lépés érvénytelen \\
        \hline
        \textbf{Tevékenység} & A metódus a snapszer szigorú lépésszabályait érvényesíti. Ha a pakli nem üres és nincs takarva, minden lépést engedélyez. Zárt vagy elfogyott pakli (végjáték) esetén sorrendben ellenőrzi a kötelező érvényű szabályokat: 1. Színkényszer (ha van a hívott színből a kézben, azonos színt kell tenni), 2. Felülütési kényszer (ha van erősebb lap az adott színből, azt kötelező megjátszani), 3. Adukényszer (ha nincs a hívott színből, de van adu a kézben, akkor kötelező adut tenni). Ha a játékos sem azonos színű lappal, sem aduval nem rendelkezik, bármilyen lapot érvényesen eldobhat. \\
        \hline
    \end{tabular}
    \caption{Az \texttt{isValidMove} metódus működésének részletezése a tényleges algoritmus alapján}
    \label{tab:actual_isvalidmove_spec}
\end{table}

\begin{table}[H]
    \centering
    \begin{tabular}{| p{0.22\textwidth} | p{0.73\textwidth} |}
        \hline
        \multicolumn{2}{|l|}{\textbf{Az \texttt{Engine} osztály \texttt{evaluateStrongerCard} metódusának megvalósítása:}} \\
        \hline
        \textbf{Bemenő adat} & \texttt{ledCard: Card} -- az ütésben elsőként kijátszott (hívott) kártyalap; \texttt{followedCard: Card} -- a másodikként megjátszott kártyalap \\
        \hline
        \textbf{Kimenő adat} & \texttt{Card} -- a két lap közül a játékszabályok értelmében erősebb, azaz az ütést elvivő kártyalap objektuma \\
        \hline
        \textbf{Tevékenység} & Meghatározza az ütés nyertesét a két kijátszott lap összehasonlításával. Ha a két kártya azonos színű, a magasabb belső értékű (rangú) lapot adja vissza. Ha a színek eltérnek, de a másodikként megjátszott lap megegyezik az aktuális aduszínnel (\texttt{trumpSuit}), akkor ez az adulap nyeri az ütést. Minden egyéb esetben (vagyis ha a második lap eltérő színű, de nem adu), az elsőként kijátszott lap marad az erősebb. \\
        \hline
    \end{tabular}
    \caption{Az \texttt{evaluateStrongerCard} metódus működésének részletezése a tényleges algoritmus alapján}
    \label{tab:actual_evaluatestrongercard_spec}
\end{table}


\begin{table}[H]
    \centering
    \begin{tabular}{| p{0.22\textwidth} | p{0.73\textwidth} |}
        \hline
        \multicolumn{2}{|l|}{\textbf{Az \texttt{Engine} osztály \texttt{switchTrumpCard} metódusának megvalósítása:}} \\
        \hline
        \textbf{Bemenő adat} & \texttt{player: Player} -- az aducserét kezdeményező játékos objektuma \\
        \hline
        \textbf{Kimenő adat} & \texttt{Object} -- egy objektum, amely tartalmaz egy \texttt{success} (boolean) tulajdonságot a művelet sikerességéről, illetve sikertelen kísérlet esetén egy magyarázó \texttt{message} (string) hibaüzenetet is tartalmaz \\
        \hline
        \textbf{Tevékenység} & A metódus az aducsere logikáját oldja meg. Először ellenőrzi, hogy a pakli nincs-e lezárva (takarva), mivel a takarás után már nem lehet adut cserélni. Utána ellenőrzi, hogy az adott játékos nem null és ő van-e soron. Ezután tovább ellenőrzi, hogy van-e az adu színnel eláttot alsó a kezében, ha ezek a feltételek igazak akkor visszatér egy objektummal, amiben egy igaz boolean szerepel, ha nem akkor hamis boolean-al eláttot objektummal tér vissza a metódus. \\
        \hline
    \end{tabular}
    \caption{A \texttt{switchTrumpCard} metódus működésének részletezése a tényleges algoritmus alapján}
    \label{tab:actual_switchtrumpcard_spec}
\end{table}

\begin{table}[H]
    \centering
    \begin{tabular}{| p{0.22\textwidth} | p{0.73\textwidth} |}
        \hline
        \multicolumn{2}{|l|}{\textbf{A \texttt{Room} osztály \texttt{broadcastState} metódusának megvalósítása:}} \\
        \hline
        \textbf{Bemenő adat} & Nem rendelkezik bemeneti paraméterrel, azonban a \texttt{Room} és az \texttt{Engine} osztályok belső állapotát használja fel az adatok kinyeréséhez. \\
        \hline
        \textbf{Kimenő adat} & Klasszikus visszatérési értéke nincs (\texttt{void}), viszont a hálózati kapcsolaton keresztül aszinkron módon egy \texttt{gameStateUpdate} eseményt bocsát ki a kliensek felé. \\
        \hline
        \textbf{Tevékenység} & Ez a metódus felel azért, hogy mindkét kliensnek szinkronizált játékfelülete legyen egyidejűleg. A rendszer végigiterál az \texttt{Engine} osztályban tárolt két játékoson, hogy mindenkinek egy perszonalizált adatcsomagot (\texttt{gameState}) generáljon. Biztonságos marad a hálózati kommunikáció, hiszen az összeállított objektumba nem kerülnek be az ellenfél rejtett kártyái, hanem az csak a játékos saját lapjait, a publikus pontszámokat, illetve a játéktér látható elemeit (adulap, pakli mérete) tartalmazza. Minden egyes eseményt követően, ami a játékfelületet befolyásolja, ez a metódus lefut és frissíti a kliensek állapotát. A végrehajtás utolsó lépéseként a generált csomag kiküldésre kerül az adott felhasználó egyedi \texttt{socketId}-jére. \\
        \hline
    \end{tabular}
    \caption{A \texttt{broadcastState} metódus működésének részletezése a tényleges algoritmus alapján}
    \label{tab:actual_broadcaststate_spec}
\end{table}

\subsection{A szerveroldali belépési pont (server.js)}
\label{subsec:server_js}

A \texttt{server.js} fájl az alkalmazás backend központi belépési pontja. Ez az állomány felelős a webszerver elindításáért, a kliens és a szerver közötti hálózati kapcsolatok felépítéséért, valamint a játék logikáját és az adatbázist összekötő hídrendszer működtetéséért. A modul a Node.js környezetben fut, és több kulcsfontosságú technológiát (Express.js \cite{expressjs}, Socket.IO \cite{socketio}, Sequelize \cite{sequelize}) integrál egyetlen egységes rendszerbe.

A modul legfontosabb felelősségi körei a következők:

\begin{itemize}
    \item \textbf{Szerver inicializálása és REST API végpontok:} Az Express.js keretrendszer segítségével itt történik a HTTP szerver konfigurálása, a CORS (Cross-Origin Resource Sharing) házirend beállítása, valamint az alapvető REST API végpontok (regisztráció, bejelentkezés, statisztikák lekérése) útvonalválasztóinak (routing) bekötése. Továbbá ez a rész felelős a lefordított kliensoldali fájlok (statikus frontend) kiszolgálásáért is.
    
    \item \textbf{Hitelesítés és jogosultságkezelés:} A játékosok csatlakozásakor a szerver JSON Web Token (JWT) segítségével azonosítja a regisztrált felhasználókat. Vendég (nem bejelentkezett) játékosok esetén ellenőrzi, hogy a választott név nem ütközik-e már létező fiókkal, illetve nem játszik-e éppen valaki ugyanazzal a névvel a szerveren, ezzel elkerülve a névütközéseket.
    
    \item \textbf{Várólista (Matchmaking) és szobakezelés:} A szerver egy memóriában tárolt várólistán (queue) tartja nyilván a meccskereső játékosokat. Amint legalább két játékos összegyűlik, a rendszer automatikusan párosítja őket, létrehoz egy egyedi azonosítóval ellátott szobát (egy \texttt{Room} példányt), és mindkét klienst ebbe a zárt Socket.IO csatornába irányítja át.
    
    \item \textbf{Valós idejű eseménykezelés (WebSockets):} A folyamatban lévő mérkőzések során a szerver folyamatosan hallgatja a kliensektől érkező eseményeket (például kártyarakás: \texttt{playerMove}, aducsere: \texttt{switchTrumpCard}, takarás: \texttt{close}), és azokat feldolgozásra a megfelelő szobapéldány játékmotorjának (\texttt{Engine}) továbbítja.
    
    \item \textbf{Kapcsolatszakadás és adatperzisztencia:} A rendszer robusztusan kezeli a váratlan hálózati megszakadásokat (\texttt{disconnect} esemény). Ha egy játékos kilépés vagy internetes hiba miatt elhagyja a folyamatban lévő meccset, a szerver a bent maradt játékost hirdeti ki győztesnek. A mérkőzések végeztével a \texttt{server.js} a Sequelize ORM segítségével aszinkron módon elmenti a meccs eredményét a \texttt{History} táblába, frissíti a résztvevők összesített statisztikáit a \texttt{Stats} táblában, végül pedig felszabadítja a lezárt szoba által lefoglalt memóriát.
\end{itemize}

\subsection{REST API végpontok dokumentációja}
\label{subsec:api_docs}

Kommunikációs hidat képeznek a REST API végpontok a frontend és a backend adatbázis között. A beérkező kérések feldolgozását és az adatáramlás irányítását a vezérlő (controller) funkciók végzik. JSON Web Token (JWT) \cite{jsonwebtoken} alapú hitelesítés védi a privát végpontokat, így a szenzitív adatokhoz (mint a profil statisztikái) csak és kizárólag a megfelelő jogosultsággal rendelkező felhasználók férhetnek hozzá.


\begin{table}[H]
    \centering
    \begin{tabular}{| p{0.12\textwidth} | p{0.28\textwidth} | p{0.5\textwidth} |}
        \hline
        \textbf{Metódus} & \textbf{Végpont (Endpoint)} & \textbf{Funkció és leírás} \\
        \hline
        \texttt{POST} & \texttt{/api/auth/register} & Új felhasználói fiókot hoz létre az adatbázisban. Ellenőrzi a jelszó és a név minimális és maximális hosszát, majd sikeres validáció után egy adatbázis-tranzakció keretében szinkronban rögzíti az \texttt{Auth} és \texttt{Player} modelleket. \\
        \hline
        \texttt{POST} & \texttt{/api/auth/login} & Hitelesíti a kliens által beküldött belépési adatokat. A \texttt{bcrypt} algoritmussal összehasonlítja a titkosított jelszavakat, majd sikeres belépés esetén egy 30 napig érvényes, aláírt JWT tokennel tér vissza. \\
        \hline
        \texttt{GET} & \texttt{/api/stats/getStats} & Lekérdezi egy adott játékos összesített profiladatait. Kizárólag érvényes JWT tokennel hívható, amelyből a szerver visszafejti a felhasználó azonosítóját, majd válaszként szolgáltatja a \texttt{Stats} és a \texttt{History} táblákból a legutóbbi tíz mérkőzés eredményét. \\
        \hline
    \end{tabular}
    \caption{A fontosabb REST API végpontok és vezérlők specifikációja}
    \label{tab:api_endpoints}
\end{table}

\subsection{Rendszeresemények naplózása (Logging)}
\label{subsec:logging}

A szerveroldali működés nyomon követése, az anomáliák kiszűrése és a hibák hatékony diagnosztizálása érdekében az alkalmazás egy egyedi, fájlalapú naplózó (logger) modult alkalmaz. A weboldal ezen komponense a Node.js beépített fájlkezelő (\texttt{fs}) és útvonalválasztó (\texttt{path}) moduljaira épül, így nem igényel harmadik féltől származó, külső függőséget.

A naplózó modul főbb jellemzői és szoftvertervezési megoldásai a következők:

\begin{itemize}
    \item \textbf{Automatikus könyvtárkezelés:} A modul inicializálásakor ellenőrzi a naplófájlok számára dedikált "logs" könyvtár meglétét az operációs rendszerben. Amennyiben a mappa nem létezik, a script (\texttt{fs.mkdirSync}) automatikusan létrehozza azt, megelőzve az Input/Output hibákat.
    
    \item \textbf{Strukturált adatrögzítés:} Az exportált \texttt{logMatchSearch} függvény egy egységes, könnyen olvasható formátumban rögzíti az eseményeket. Minden bejegyzés tartalmazza a művelet pontos, lokalizált időbélyegét, a játékos nevét, valamint a csatlakozási kísérlet státuszát.
    
    \item \textbf{Hibakövetés:} A modul megkülönbözteti a sikeres és sikertelen meccskeresési eseményeket. Hibás kísérletek esetén (például érvénytelen JWT token, foglalt név, vagy hibás adatbázis-kapcsolat) a naplósor kiegészül a hiba pontos okával is, ami elengedhetetlen a későbbi rendszerüzemeltetési feladatokhoz.
\end{itemize}

\subsection{Alkalmazott szoftverkönyvtárak és függőségek}
\label{subsec:dependencies}

A projekt fejlesztése során a \texttt{package.json} fájlban meghatározott, célirányosan kiválasztott külső könyvtárak (függőségek) kerültek felhasználásra. A beépülő al-könyvtárak (tranzitív függőségek) mellőzésével a rendszer legfőbb építőelemei és azok felelősségi körei a következők:

\begin{itemize}
    \item \textbf{Express (\texttt{express} v5.2.1):} A Node.js legelterjedtebb minimalista webes keretrendszere, amely a REST API végpontok kialakításáért, a HTTP kérések (útvonalválasztás és middleware-ek) kezeléséért, valamint a statikus frontend fájlok kiszolgálásáért felel \cite{expressjs}.
    
    \item \textbf{Socket.IO (\texttt{socket.io} v4.8.3):} A valós idejű, kétirányú hálózati kommunikációt (WebSockets alapokon) biztosító keretrendszer, amely a folyamatos és gyors adatcserét, ezáltal a tényleges játékmenet szinkronizálását végzi a kliensek és a szerver között \cite{socketio}.
    
    \item \textbf{Sequelize (\texttt{sequelize} v6.37.8):} Egy fejlett, ígéret-alapú (promise-based) ORM (Object-Relational Mapper) eszköz, amely az adatbázis-műveleteket absztrakt, objektumorientált módon kezeli, kiváltva a nyers SQL lekérdezések manuális írását \cite{sequelize}.
    
    \item \textbf{PostgreSQL csomagok (\texttt{pg} v8.20.0, \texttt{pg-hstore} v2.3.4):} A Sequelize ORM számára elengedhetetlen natív adatbázis-meghajtók, amelyek a Node.js és a PostgreSQL adatbázis-szerver közötti alacsony szintű kommunikációt valósítják meg \cite{node-postgres}.  
    
    \item \textbf{JSON Web Token (\texttt{jsonwebtoken} v9.0.3):} A token alapú, állapotmentes (stateless) felhasználói hitelesítést (autentikációt) és jogosultságkezelést biztosító könyvtár, amely titkosított formában tárolja a munkamenet adatait \cite{jsonwebtoken}.
    
    \item \textbf{Bcrypt.js (\texttt{bcryptjs} v3.0.3):} A felhasználói jelszavak biztonságos, egyirányú titkosításáért (hashing) felelős modul, amellyel garantálható, hogy az adatbázis kompromittálódása esetén sem nyerhetőek ki a jelszavak nyílt szövegként \cite{bcryptjs}.
    
    \item \textbf{CORS (\texttt{cors} v2.8.6):} Biztonsági köztesréteg (middleware) az Express-hez, amely a Cross-Origin Resource Sharing házirend konfigurálásával engedélyezi az eltérő tartományokból érkező böngészős kéréseket \cite{cors-pkg}.
    
    \item \textbf{Dotenv (\texttt{dotenv} v17.2.3):} A környezeti változók (például adatbázis-kapcsolati azonosítók, titkosító kulcsok, szerver portok) egy elkülönített \texttt{.env} fájlból történő, biztonságos betöltéséért felel \cite{dotenv}.
    
    \item \textbf{Tailwind CSS (\texttt{tailwindcss} v4.1.18, \texttt{@tailwindcss/vite}):} A kliensoldali grafikus felület (UI) atomi osztályokon alapuló, reszponzív formázását végző CSS keretrendszer és annak build-eszköz (Vite) integrációja \cite{tailwindcss}.
    
    \item \textbf{Nodemon (\texttt{nodemon} v3.1.14):} Fejlesztési célú segédeszköz (devDependency), amely a forráskód mentésekor automatikusan újraindítja a Node.js háttérszervert, ezzel jelentősen gyorsítva az iterációs folyamatot \cite{nodemon}.
\end{itemize}


\section{Élesítési környezet és infrastruktúra}
\label{sec:deployment}

A weboldal teljes egésze modern, felhőalapú infrastruktúrán fut, így mind a \url{https://snapszer.com} oldalt kiszolgáló szerver, mind a háttérben futó adatbázis a Microsoft Azure \cite{azure} környezetében kapott helyet.

Az adattárolásért egy különálló, bérelt PostgreSQL adatbázis-szerver (Azure Database for PostgreSQL) felel, amely a regisztrációkat, a játékosok pontszámait és a mérkőzések részletes történetét rögzíti. Azzal, hogy a weboldal és az adatbázis is a Microsoft felhőszolgáltatását használja, a rendszer rendkívül stabil és gyors maradhat.

A felhőalapú működésnek köszönhetően a játékszerver minden erőforrását a mérkőzések zökkenőmentes lebonyolítására tudja fordítani. Emellett a Microsoft Azure biztosítja a folyamatos elérhetőséget, a rendszer könnyű bővíthetőségét. A biztonságot tovább fokozza az Azure beépített hálózati tűzfala, a DDoS-támadások elleni automatikus védelmi rendszere, valamint az adatforgalom kötelező alapú titkosítása (TLS/SSL), amelyek együttesen proaktív módon védik az alkalmazást a külső kiberfenyegetésekkel szemben.


\section{Rendszertesztelés és a megvalósult használati esetek}
\label{sec:testing_use_cases}

A fejlesztési fázis lezárultával a szoftver minőségbiztosítása és az üzleti logika helyességének ellenőrzése következett. Ennek alapját a Behaviour-Driven Development (BDD) módszertan adta, amely a rendszer működését a felhasználói viselkedés felől közelíti meg. 

Az alábbi táblázat az alkalmazás legfontosabb használati eseteit (Use Case) és azok tesztelésének eredményeit mutatja be a szabványos \textit{Given-When-Then} (Adott-Amikor-Akkor) formátumban. A jobb oldali oszlopban található „OK” minősítés igazolja, hogy a végleges implementáció a tesztek során az elvárt viselkedést produkálta, így az aktuális megvalósítás fedi az előzetes terveket.

\vspace{1em}

\begin{longtable}{| c | m{3cm} | c | m{6cm} | c |}
    \caption{A rendszer használati eseteinek részletezése és teszteredményei}
    \label{tab:use-case-scenarios} \\
    
    \hline
    \textbf{\#} & \textbf{Eset} & \multicolumn{2}{c|}{\textbf{Leírás}} & \textbf{Eredmény} \\
    \hline
    \endfirsthead
    
    \hline
    \textbf{\#} & \textbf{Eset} & \multicolumn{2}{c|}{\textbf{Leírás}} & \textbf{Eredmény} \\
    \hline
    \endhead
    
    \hline
    \endlastfoot

    \multirow{3}{*}{1.} & \multirow{3}{3cm}{Sikeres regisztráció} 
    & GIVEN & Vendégként tartózkodom a regisztrációs oldalon. & \cellcolor{green!40} \\* \cline{3-4}
    & & WHEN & Kitöltöm a szükséges adatokat (felhasználónév, jelszó) és rákattintok a „Regisztráció” gombra. & \cellcolor{green!40}\textbf{OK} \\* \cline{3-4}
    & & THEN & A rendszer létrehozza a fiókomat az adatbázisban. & \cellcolor{green!40} \\ \hline

    \newpage

    \multirow{3}{*}{2.} & \multirow{3}{3cm}{Sikeres bejelentkezés} 
    & GIVEN & Vendégként tartózkodom a bejelentkezési oldalon. & \cellcolor{green!40} \\* \cline{3-4}
    & & WHEN & Megadom a regisztrált felhasználónevemet és jelszavamat, majd a „Belépés” gombra kattintok. & \cellcolor{green!40}\textbf{OK} \\* \cline{3-4}
    & & THEN & A rendszer hitelesíti az adatokat, generálja a tokent, és a főmenübe irányít. & \cellcolor{green!40} \\ \hline

    \multirow{3}{*}{3.} & \multirow{3}{3cm}{Játék keresése (Matchmaking)} 
    & GIVEN & Játékosként be vagyok jelentkezve, és a főmenüben vagyok. & \cellcolor{green!40} \\* \cline{3-4}
    & & WHEN & Rákattintok a játék keresése gombra. & \cellcolor{green!40}\textbf{OK} \\* \cline{3-4}
    & & THEN & A rendszer felvesz a várólistára, majd amint van elérhető ellenfél, létrehoz egy közös szobát. & \cellcolor{green!40} \\ \hline

    \multirow{3}{*}{4.} & \multirow{3}{3cm}{Játék keresése egy foglalt névvel} 
    & GIVEN & Vendégként vagyok a főmenüben. & \cellcolor{green!40} \\* \cline{3-4}
    & & WHEN & Beírom a becenevemet a keresőbe és rákattintok a játék keresése gombra. & \cellcolor{green!40}\textbf{OK} \\* \cline{3-4}
    & & THEN & A rendszer észreveszi, hogy már játszik egy ilyen nevű játékos és jelzni nekem a főmenűben, hogy foglalt a név. & \cellcolor{green!40} \\ \hline

    \multirow{3}{*}{5.} & \multirow{3}{3cm}{Érvényes kártya kijátszása} 
    & GIVEN & Játékosként egy mérkőzésben vagyok, és én következem. & \cellcolor{green!40} \\* \cline{3-4}
    & & WHEN & A kezemben lévő kártyák közül az ütéskényszer vagy színkényszer szabályainak megfelelően kijátszok egyet. & \cellcolor{green!40}\textbf{OK} \\* \cline{3-4}
    & & THEN & A rendszer elfogadja a lépést, a lap megjelenik az asztalon, és a sor az ellenfélre száll. & \cellcolor{green!40} \\ \hline

    \multirow{3}{*}{6.} & \multirow{3}{3cm}{Érvénytelen kártya kijátszása} 
    & GIVEN & Játékosként egy mérkőzésben vagyok, és én következem. & \cellcolor{green!40} \\* \cline{3-4}
    & & WHEN & A kezemben lévő kártyák közül az ütéskényszer vagy színkényszer szabályainak nem megfelelően kijátszok egyet. & \cellcolor{green!40}\textbf{OK} \\* \cline{3-4}
    & & THEN & A rendszer elutasítja a lépést, a lap egy vastag piros keretet kap és rezeg egy rövid ideig jelezve, hogy a lépésem illegális. & \cellcolor{green!40} \\ \hline

    \multirow{3}{*}{7.} & \multirow{3}{3cm}{Bemondás (20/40)} 
    & GIVEN & Hívó félként van egy királyom és egy felsőm azonos színből. & \cellcolor{green!40} \\* \cline{3-4}
    & & WHEN & Kijátszom a 20-at vagy 40-et érő páros egyik tagját (királyt vagy felsőt) a gombra kattintva. & \cellcolor{green!40}\textbf{OK} \\* \cline{3-4}
    & & THEN & A rendszer érvényesíti a bemondást, jóváírja a plusz pontokat, és megjeleníti a lapot. & \cellcolor{green!40} \\ \hline

    \multirow{3}{*}{8.} & \multirow{3}{3cm}{Ütés kiértékelése} 
    & GIVEN & Mindkét játékos kijátszotta a kártyáját az adott körben. & \cellcolor{green!40} \\* \cline{3-4}
    & & WHEN & A rendszer megvizsgálja a lapokat az aduszín és a laperősség alapján. & \cellcolor{green!40}\textbf{OK} \\* \cline{3-4}
    & & THEN & A rendszer megállapítja az ütés nyertesét, jóváírja a pontokat, és frissíti a klienseket. & \cellcolor{green!40} \\ \hline

    \multirow{3}{*}{9.} & \multirow{3}{3cm}{Aducsere} 
    & GIVEN & Játékosként én következem, és a kezemben van az aktuális aduszínhez tartozó alsó kártyalap. & \cellcolor{green!40} \\* \cline{3-4}
    & & WHEN & Rákattintok a játéktéren elhelyezett adulapra. & \cellcolor{green!40}\textbf{OK} \\* \cline{3-4}
    & & THEN & A rendszer kicseréli az adulapot a kezemben lévő alsóra, és frissíti a játéktér állapotát mindkét kliensnél. & \cellcolor{green!40} \\ \hline

    \multirow{3}{*}{10.} & \multirow{3}{3cm}{Helytelen aducsere} 
    & GIVEN & Játékosként én következem, és a kezemben nincs az aktuális aduszínhez tartozó alsó kártyalap. & \cellcolor{green!40} \\* \cline{3-4}
    & & WHEN & Rákattintok a játéktéren elhelyezett adulapra. & \cellcolor{green!40}\textbf{OK} \\* \cline{3-4}
    & & THEN & A rendszer elutasítja az eseményt és az adulap egy vastag piros keretet kap és rezeg egy rövid ideig. & \cellcolor{green!40} \\ \hline

    \multirow{3}{*}{11.} & \multirow{3}{3cm}{Pakli takarása (Zárás)} 
    & GIVEN & Játékosként én vagyok soron, és a húzópakliban (talonban) még legalább 3 kártyalap található. & \cellcolor{green!40} \\* \cline{3-4}
    & & WHEN & Egy gombnyomással kezdeményezem a pakli takarását. & \cellcolor{green!40}\textbf{OK} \\* \cline{3-4}
    & & THEN & A rendszer lezárja a paklit, vizuálisan jelzi ezt a talonon, és egy felugró ablakban értesíti a résztvevőket. & \cellcolor{green!40} \\ \hline

    \multirow{3}{*}{12.} & \multirow{3}{3cm}{Kör megnyerése (66 pont)} 
    & GIVEN & Egy mérkőzés során véget ér egy ütés. & \cellcolor{green!40} \\* \cline{3-4}
    & & WHEN & Az ütést elvivő játékos összegyűjtött pontjainak száma eléri a 66-ot. & \cellcolor{green!40}\textbf{OK} \\* \cline{3-4}
    & & THEN & A rendszer lezárja a kört, kihirdeti a nyertest, és kiosztja a megfelelő meccspontokat. & \cellcolor{green!40} \\ \hline

    \multirow{3}{*}{13.} & \multirow{3}{3cm}{Kör elvesztése (66 pont az ellenfélnek)} 
    & GIVEN & Egy mérkőzés során véget ér egy ütés. & \cellcolor{green!40} \\* \cline{3-4}
    & & WHEN & Az ellenfelem az általa elvitt ütésekkel összegyűjt legalább 66 pontot. & \cellcolor{green!40}\textbf{OK} \\* \cline{3-4}
    & & THEN & A rendszer lezárja a kört, az ellenfelet hirdeti ki a kör nyertesének, és neki oszt meccspontot. & \cellcolor{green!40} \\ \hline

    \multirow{3}{*}{14.} & \multirow{3}{3cm}{Meccs megnyerése (7 meccspont)} 
    & GIVEN & Egy aktív mérkőzés során a rendszer kiértékel egy újonnan megnyert kört. & \cellcolor{green!40} \\* \cline{3-4}
    & & WHEN & A frissen megnyert körért kapott pontokkal az én összesített meccspontjaim száma eléri a 7-et. & \cellcolor{green!40}\textbf{OK} \\* \cline{3-4}
    & & THEN & A rendszer lezárja a mérkőzést, kihirdet engem győztesnek, és elmenti az eredményt az adatbázisba. & \cellcolor{green!40} \\ \hline

    \multirow{3}{*}{15.} & \multirow{3}{3cm}{Meccs elvesztése (7 meccspont az ellenfélnek)} 
    & GIVEN & Egy aktív mérkőzés során a rendszer kiértékel egy elvesztett kört. & \cellcolor{green!40} \\* \cline{3-4}
    & & WHEN & Az ellenfelem megnyeri a kört, és az azért kapott pontokkal az ő meccspontjainak száma eléri a 7-et. & \cellcolor{green!40}\textbf{OK} \\* \cline{3-4}
    & & THEN & A rendszer a mérkőzést lezárja, az ellenfelet hirdeti ki győztesnek, engem pedig a főmenübe irányít. & \cellcolor{green!40} \\ \hline

    \multirow{3}{*}{16.} & \multirow{3}{3cm}{Játék elhagyása (Disconnect)} 
    & GIVEN & Egy aktív mérkőzésen veszek részt. & \cellcolor{green!40} \\* \cline{3-4}
    & & WHEN & Kilépek az oldalról vagy bezárom a böngészőt. & \cellcolor{green!40}\textbf{OK} \\* \cline{3-4}
    & & THEN & A rendszer érzékeli a kilépést, a bent maradt játékost hirdeti ki győztesnek, és lezárja a szobát. & \cellcolor{green!40} \\ \hline

    \multirow{3}{*}{17.} & \multirow{3}{3cm}{Statisztikák megtekintése} 
    & GIVEN & Játékosként bejelentkezve a főmenüben tartózkodom. & \cellcolor{green!40} \\* \cline{3-4}
    & & WHEN & Rákattintok a statisztikák gombra. & \cellcolor{green!40}\textbf{OK} \\* \cline{3-4}
    & & THEN & A rendszer lekéri és listázza az eddig elért eredményeimet (győzelmek, meccsek száma). & \cellcolor{green!40} \\ \hline

    \multirow{3}{*}{18.} & \multirow{3}{3cm}{Játék keresés megszakítása} 
    & GIVEN & A főmenüben tartózkodom és a játék keresése zajlik. & \cellcolor{green!40} \\* \cline{3-4}
    & & WHEN & Rákattintok az "X" gombra a játék keresésnél. & \cellcolor{green!40}\textbf{OK} \\* \cline{3-4}
    & & THEN & A játék keresése megszakad és újra megjelenik a játék keresése gomb. & \cellcolor{green!40} \\ \hline

\end{longtable}

\subsection{Tesztelés által befolyásolt döntések}
\label{subsec:test_driven_decisions}

A szoftver fejlesztése során a folyamatos manuális tesztelés több olyan hiányosságra is rávilágított, amelyek nem feltétlenül programhibák voltak, hanem főleg a felhasználói élményt és a játékmenet érthetőségét rontották. Ezen visszajelzések alapján több kritikus módosítást is eszközöltem a rendszeren:

\begin{itemize}
    \item \textbf{Vizuális visszacsatolás a játékmenet eseményeiről:} 
    A kezdeti tesztfázisban a pakli takarása (zárása) vagy a 20/40 pontos bemondások során az ellenfél nem kapott egyértelmű jelzést arról, hogy mi történt, így a játékmenet zavarossá vált. Ezt orvosolva bevezettem egy értesítési rendszert, amely azonnal tájékoztatja a másik felet ezekről a stratégiai lépésekről, biztosítva a játék folyamatos követhetőségét. Ezeket az értesítéseket úgynevezett felugró ablakokkal oldottam meg, amellyek színnel is adnak információt: pl. narancsal jelenik meg, amikor az ellenfél játszik ki 20/40-et és kékkel ha mi játszunk ki.

    \item \textbf{A körök kiértékelésének ütemezése:} 
    Eredetileg, amikor a második játékos is kijátszotta a kártyáját, a rendszer azonnal elvégezte a kiértékelést, és a lapok azonnal eltűntek az asztalról. Ez zavaró volt, mivel a felhasználóknak nem maradt idejük vizuálisan feldolgozni az ütést. Ennek javítására bevezettem egy 2 másodperces késleltetést, amely alatt mindkét kijátszott kártya látható marad. Ez a rövid szünet segít a játékosoknak követni a pontszámítást és az ütés irányát, mielőtt a rendszer továbblépne a következő körre.

    \item \textbf{A várólistáról való kilépés lehetősége:} 
    A hálózati kommunikáció tesztelésekor kiderült, hogy miután a játékos elindította a keresést, nem volt lehetősége meggondolnia magát és elhagyni a várólistát. Ez kényelmetlen helyzetet teremtett hosszú várakozási idő esetén. A probléma megoldására a kereső felületet kiegészítettem egy "X" (megszakítás) gombbal, amely lehetővé teszi a keresési folyamat biztonságos leállítását és a főmenübe való visszatérést anélkül, hogy az oldalt újra kellene tölteni.
\end{itemize}

Ezek a módosítások jelentősen növelték a szoftver ergonómiáját és a játékélmény folytonosságát, bizonyítva a tesztelés fontosságát.